\documentclass[11pt,a4paper]{article}
\usepackage{DDTA_METHODOLOGY_GUIDE_STYLE_R1}
\usepackage{amsmath,amssymb}

\providecommand{\BA}{\textit{Base Analysis}}
\providecommand{\BAR}{\textit{BAReferent}}
\providecommand{\BAP}{\textit{BAProposition}}
\providecommand{\stabletag}{\ddtatag{DDTAGreen}{CURRENT / PRESERVED}}
\providecommand{\Rtwentyfivetag}{\ddtatag{DDTATeal}{R25 EVIDENCE}}
\providecommand{\candidatetag}{\ddtatag{DDTAAmber}{CANDIDATE / NOT ADMITTED}}
\providecommand{\rejectedtag}{\ddtatag{DDTARed}{REJECTED}}
\providecommand{\pendingexampletag}{\ddtatag{DDTABlue}{CORE VIEW / NO EXAMPLES}}

\begin{document}
\selectlanguage{italian}
\hypersetup{pdftitle={DDTA Base Analysis Core Guide - Candidate R1}}
\fancyhead[L]{\sffamily\small\color{DDTAGray}DDTA - Base Analysis Core Guide}
\fancyhead[R]{\sffamily\small\color{DDTAAmber}CORE GUIDE - CANDIDATE R1 - R42}
\fancyfoot[L]{\sffamily\scriptsize\color{DDTAGray}Core semantic/syntax guide - no worked examples and no discrimination question pack}

\begin{titlepage}
\thispagestyle{empty}
\vspace*{8mm}
{\sffamily\bfseries\color{DDTANavy}\fontsize{15}{18}\selectfont Documentation-Driven Threat Analysis}\par
\vspace{5mm}
{\sffamily\bfseries\color{DDTANavy}\fontsize{29}{33}\selectfont Base Analysis Core Guide}\par
\vspace{2mm}
{\sffamily\color{DDTABlue}\fontsize{18}{22}\selectfont Candidate R1 - Semantic and Syntax Core}\par

\vspace{9mm}
\begin{tcolorbox}[enhanced,arc=3pt,boxrule=0pt,colback=DDTANavy,left=12pt,right=12pt,top=11pt,bottom=11pt]
{\color{white}\sffamily\bfseries\large Base semantica e sintattica della BA, volutamente senza worked examples e senza question pack}\par
\vspace{2mm}
{\color{white!88}\small Questa guida isola il core semantico e sintattico della BA: definizioni, firme, roleKey, cardinalita', condition language, strutture locali, provenance, projection, candidate construct e open register. Worked examples e question-driven discrimination sono intenzionalmente esclusi e sono forniti dalla Complete Guide companion nello stesso checkpoint R42.}
\end{tcolorbox}

\vspace{8mm}
\begin{tabularx}{\textwidth}{L{47mm}Y}
\toprule
\textbf{Stato} & \texttt{CANDIDATE / NON-NORMATIVE / NOT PROMOTION ELIGIBLE} \\
\textbf{Scopo di questa revisione} & Completezza semantica e sintattica del core BA, separata dai layer pedagogici e interrogativi. \\
\textbf{Authority BA corrente} & \texttt{DDTA\_BASE\_ANALYSIS\_OPERATIONAL\_GUIDE\_R3.tex} + BA0 R1 / BA1 R1 / BA2 R3 / BA3 R1 / BA4 R1 / BA5 R1. \\
\textbf{Operator basis corrente} & 14 top-level operators. Nessun quindicesimo operatore e' ammesso. \\
\textbf{Repository baseline di partenza} & \texttt{bea1e6567d3f861294378b201648b0853b3c8a43} \\
\textbf{Layer esempi} & Intenzionalmente assente; disponibile nella Complete Guide companion di R42. \\
\textbf{Data} & 11 settembre 2026 \\
\bottomrule
\end{tabularx}

\vfill
\begin{ddtaopen}[title={Promotion stop esplicito}]
Questa Core Guide non e' un successore autonomo della R3: e' una vista intenzionalmente ridotta al contratto semantico/sintattico. La Complete Guide companion conserva questo core e aggiunge worked examples e question-driven discrimination.
\end{ddtaopen}
\end{titlepage}

\tableofcontents
\clearpage

\section{Scopo, lineage e regola di costruzione}

Questa guida separa deliberatamente due obiettivi che nelle revisioni precedenti erano intrecciati:
\begin{enumerate}
  \item \textbf{completezza semantica e sintattica}: ogni costrutto deve avere significato, firma, ruoli, cardinalita', confini e stato;
  \item \textbf{completezza pedagogica}: ogni costrutto deve essere accompagnato da esempi source-grounded e controesempi.
\end{enumerate}

La R2 chiude soltanto il primo obiettivo. Il secondo e' esplicitamente rinviato alla revisione successiva.

\begin{ddtacore}[title={Regola cumulativa}]
Una candidate successore non puo' perdere una distinzione stabile della R3 per essere piu' breve. Ogni contenuto R3 viene quindi classificato come preservato, espanso, qualificato oppure esplicitamente rinviato al layer di esempi. Nessuna evidence R25 diventa authority soltanto perche' e' inclusa in questa guida.
\end{ddtacore}

\subsection{Legenda epistemica}
\begin{tabularx}{\textwidth}{L{55mm}Y}
\toprule
\textbf{Etichetta} & \textbf{Significato} \\
\midrule
\stabletag & Regola gia' authority corrente in R3/BA0--BA5. \\
\Rtwentyfivetag & Evidenza positiva o pressure confermata dal secondo ciclo DermaTriage. \\
\candidatetag & Costrutto sopravvissuto a test parziali ma non formalmente ammesso. \\
\rejectedtag & Working hypothesis respinta nello scope corrente. \\
\opentag & Questione semanticamente reale ancora non chiusa. \\
\pendingexampletag & La semantica e' descritta; esempi e question pack sono intenzionalmente esclusi da questa vista core. \\
\bottomrule
\end{tabularx}

\section{Che cos'e' la Base Analysis}
\stabletag

Per una baseline documentale governata, la \BA{} e' la rappresentazione analitica condivisa e methodology-neutral del significato di progetto necessario ai consumer downstream DDTA.

\begin{lstlisting}
governed project documentation
    = project authority

accepted Base Analysis
    = shared analytical baseline

projection / threat analysis / tests
    = downstream consumers
\end{lstlisting}

La BA puo' normalizzare, rendere riusabile e rendere interrogabile un significato gia' governato. Non puo' completare silenziosamente la documentazione, inventare un commitment, scegliere una interpretazione non supportata o trasformare una domanda downstream in project truth.

\begin{ddtacore}[title={Authority direction}]
Documentazione governata $\rightarrow$ BA $\rightarrow$ projection/analysis. Il feedback puo' risalire come richiesta di review, ma una modifica di project truth avviene soltanto tramite una nuova decisione documentale governata e successiva revalidation della BA.
\end{ddtacore}

\section{Contratto BA0--BA5 e gate BA6}
\stabletag

\begin{longtable}{L{18mm}L{123mm}}
\toprule
\textbf{Layer} & \textbf{Responsabilita'} \\
\midrule\endhead
BA0 & authority boundary, minimum sufficient analysis, unresolved state, feedback direction e anti-inference. \\
BA1 & due sole famiglie first-class: \BAR{} e \BAP{}. \\
BA2 R3 & proposition semantics, top-level operators, operator-scoped roleKey, cardinalita' e strutture locali. \\
BA3 & provenance, origin state, derivation, review/freshness, continuity e revalidation. \\
BA4 & projection, traceability, interpretation e coverage downstream. \\
BA5 & canonical semantic registries, domain-scoped resolution e controlled authoring. \\
BA6 & integrated validation gate. Non viene dichiarato chiuso da questa guida. \\
\bottomrule
\end{longtable}

\begin{ddtacore}[title={BA2 authority corrente}]
Le firme di operatori, roleKey e cardinalita' normative derivano da \texttt{BA2\_RELATION\_ACTION\_VOCABULARY\_R3}. Snapshot storici incorporati in altri contratti non sostituiscono BA2 R3.
\end{ddtacore}

\section{Procedura operativa end-to-end}
\stabletag

\begin{enumerate}
  \item Pin di baseline, authority e source revision.
  \item Lettura source-first del governed source set.
  \item Identificazione del minimum shared meaning necessario.
  \item Identificazione dei significati che richiedono identita' riusabile.
  \item Creazione dei BAReferent necessari.
  \item Formulazione delle BAProposition minime usando BA2 R3.
  \item Verifica di firma, roleKey, cardinalita' e scope delle proposition.
  \item Collegamento BA3 a source evidence, origin state e revalidation context.
  \item Registrazione esplicita di ambiguity, conflict, diagnosis e \texttt{NOT SPECIFIED}.
  \item Minimum-BA acceptance per lo scope dichiarato.
  \item Regression completa verso le fonti e forbidden-inference review.
  \item Method-pressure review soltanto su counterexample concreti.
  \item Materializzazione della accepted case BA.
  \item Projection/downstream consumption secondo BA4.
  \item Governed change e BA3 revalidation quando necessario.
  \item Integrated BA6 verdict soltanto dopo l'intero protocollo previsto.
\end{enumerate}

\section{BAReferent}
\stabletag

Un \BAR{} e' un significato di progetto methodology-neutral che riceve identita' analitica indipendente quando tale identita' e' materialmente necessaria.

\subsection{Quando creare un BAReferent}
Creare un BAReferent quando il significato deve almeno una volta essere:
\begin{itemize}
  \item riusato da piu' proposition;
  \item qualificato o vincolato;
  \item correlato con un altro significato;
  \item referenziato come target distinto;
  \item confrontato tra baseline;
  \item target di provenance/revalidation indipendente;
  \item oggetto di una domanda downstream materialmente distinta.
\end{itemize}

\subsection{Quando NON creare un BAReferent}
Non creare un referent soltanto perche' compare un sostantivo nel testo, perche' un tool desidera un nodo, perche' ogni verbo dovrebbe avere un oggetto o perche' una projection grafica risulterebbe piu' comoda.

\begin{ddtacore}[title={Identity test}]
Se il significato scompare come identita' autonoma senza perdere alcuna distinzione governata, alcuna possibilita' di qualificazione o alcuna tracciabilita' necessaria, probabilmente non richiede un BAReferent separato.
\end{ddtacore}

\section{BAProposition}
\stabletag

Una \BAP{} e' una asserzione analitica formalizzata su uno o piu' BAReferent. La proposition applica un operatore e assegna BAReferent ai roleKey ammessi da quell'operatore.

\begin{lstlisting}
BAProposition
    operator -> canonical BA2 operator key
    roles    -> operator-scoped role assignments
    local    -> only operator-local structures when admitted
\end{lstlisting}

\subsection{Invarianti}
\begin{itemize}
  \item una proposition non e' il referent che descrive;
  \item roleKey sono scoped all'operatore, non proprieta' universali del referent;
  \item un ruolo obbligatorio mancante non viene riempito per default;
  \item un ruolo opzionale si usa soltanto se il significato e' governato;
  \item una proposition non deve contenere semantica non ammessa dalla propria firma;
  \item piu' proposition possono essere necessarie per preservare distinzioni diverse contenute nella stessa clausola sorgente.
\end{itemize}

\section{Criterio di dettaglio, minimalita' e stopping rule}
\stabletag

\begin{lstlisting}
more detail
    only if needed to:
      preserve a governed distinction
      answer a material downstream question
      expose a missing governed fact
\end{lstlisting}

\begin{ddtacheck}[title={Stopping test}]
Abbiamo abbastanza semantica per formulare la domanda rilevante e distinguere evidence supportata, contraddetta o \texttt{NOT SPECIFIED}? Se SI, fermarsi. Se NO, identificare la smallest missing distinction o produrre una diagnosis. Non inventare il dettaglio mancante.
\end{ddtacheck}

\section{Operator selection procedure}
\stabletag

La scelta dell'operatore avviene dopo la ricostruzione del significato, non prima.

\begin{enumerate}
  \item isolare la smallest governed semantic fact;
  \item identificare i BAReferent necessari;
  \item formulare le candidate interpretations materialmente differenti;
  \item provare prima gli operatori esistenti;
  \item verificare firma e roleKey senza adattamenti opportunistici;
  \item applicare delete test e merge-first test;
  \item verificare i confini con gli operatori semanticamente vicini;
  \item se nessun costrutto preserva il significato, registrare pressure/diagnosis invece di inventare un nuovo operatore locale;
  \item un nuovo operatore richiede formal admission separata.
\end{enumerate}

\begin{ddtaanti}[title={Verb-to-operator mapping}]
Il verbo naturale nella documentazione non determina automaticamente l'operatore BA. Parole come ``acquire'', ``invoke'', ``store'', ``select'', ``execute'', ``trigger'' e ``expose'' possono descrivere semantiche differenti a seconda del significato governato.
\end{ddtaanti}

\section{RoleKey: significato e disciplina}
\stabletag

I roleKey definiscono la funzione svolta da un BAReferent \emph{all'interno di una specifica proposition}. Non classificano universalmente il referent.

\begin{longtable}{L{45mm}L{97mm}}
\toprule
\textbf{roleKey} & \textbf{Significato operativo} \\
\midrule\endhead
\texttt{behavior} & identita' riusabile dello specifico transfer quando il comportamento di conveyance deve essere qualificato o target di altre proposition. \\
\texttt{source} & origine del content trasferito. \\
\texttt{destination} & destinatario del content trasferito. \\
\texttt{content} & significato oggetto della conveyance. \\
\texttt{actor} & soggetto/capability che compie l'azione quando la firma lo richiede o lo ammette. \\
\texttt{result} & risultato reso disponibile o risultato governato dell'operazione. \\
\texttt{input} & significato effettivamente usato dall'operazione; non ogni dato disponibile e' automaticamente input. \\
\texttt{created} & nuovo item/event project-semantic stabilito da \texttt{create}. \\
\texttt{observed} & significato esistente letto o ispezionato da \texttt{observe}. \\
\texttt{subject} & identita' la cui state/lifecycle identity e' interessata da \texttt{transition}, oppure subject locale di \texttt{satisfies}. \\
\texttt{fromState} & stato precedente governato, se espresso dalla fonte. \\
\texttt{toState} & stato risultante governato richiesto da \texttt{transition}. \\
\texttt{correlatedItem} & elemento che deve mantenere same-context binding. \\
\texttt{correlationContext} & identita' che definisce il contesto condiviso della correlation. \\
\texttt{referencer} & elemento che mantiene un riferimento direzionale verso altro significato. \\
\texttt{referenced} & target del riferimento. \\
\texttt{dependent} & significato che richiede uno o piu' prerequisiti. \\
\texttt{prerequisite} & significato necessario al dependent. \\
\texttt{consumer} & utilizzatore effettivo di un service/capability. \\
\texttt{service} & service/capability effettivamente consumato. \\
\texttt{provider} & provider del service, soltanto se governato; non e' implicito nel servizio. \\
\texttt{abstract} & significato astratto indipendentemente governato da materializzare. \\
\texttt{realization} & significato concreto che materializza l'abstract. \\
\texttt{responsibleParty} & parte cui viene assegnata responsibility/authority. \\
\texttt{responsibilityScope} & ambito/oggetto della responsibility assegnata. \\
\texttt{responsibilityKind} & tipo governato di responsibility/authority. \\
\texttt{constraintTarget} & significato cui si applica una restriction governata. \\
\texttt{constraintValue} & valore, dominio o struttura di restriction governata. \\
\texttt{classifiedReferent} & elemento cui viene assegnato un semantic kind. \\
\texttt{semanticKind} & categoria methodology-neutral governata. \\
\bottomrule
\end{longtable}

\section{Matrice normativa delle cardinalita'}
\stabletag

\begin{longtable}{L{42mm}L{65mm}L{33mm}}
\toprule
\textbf{Operator} & \textbf{Required} & \textbf{Optional / local} \\
\midrule\endhead
transfer & source 1; destination 1..*; content 1..* & behavior 0..1 \\
produce & actor 1; result 1..* & input 0..* \\
create & actor 1; created 1..* & none \\
observe & actor 1; observed 1..* & result 0..* \\
transition & subject 1; toState 1 & actor 0..1; fromState 0..1 \\
correlate & correlatedItem 1..*; correlationContext 1 & none \\
reference & referencer 1; referenced 1..* & none \\
dependOn & dependent 1; prerequisite 1..* & none \\
consumeService & consumer 1..*; service 1 & provider 0..1 \\
realize & abstract 1; realization 1..* & none \\
assignResponsibility & responsibleParty 1..*; responsibilityScope 1; responsibilityKind 1 & none \\
constrain & constraintTarget 1; constraintValue 1..* & none \\
classify & classifiedReferent 1; semanticKind 1..* & none \\
decisionRule & actor 1; input 1..*; result 1 & operator-local \texttt{rule} \\
\bottomrule
\end{longtable}

% -----------------------------------------------------------------------------
\section{Operator: transfer}
\stabletag\quad\Rtwentyfivetag\quad\pendingexampletag

\subsection{Definizione}
\texttt{transfer} rappresenta la conveyance governata di uno o piu' content da una source a una o piu' destination. Non implica di per se' protocollo, successo, ownership, storage, processing o service consumption.

\subsection{Firma canonica}
\begin{lstlisting}
transfer
    behavior    -> BAReferent [0..1]
    source      -> BAReferent [1]
    destination -> BAReferent [1..*]
    content     -> BAReferent [1..*]
\end{lstlisting}

\subsection{Campi}
\begin{tabularx}{\textwidth}{L{36mm}L{24mm}Y}
\toprule
\textbf{Role} & \textbf{Card.} & \textbf{Semantica} \\
\midrule
behavior & 0..1 & identita' del comportamento di transfer, solo se deve essere riusato/qualificato. \\
source & 1 & origine governata della conveyance. \\
destination & 1..* & uno o piu' destinatari governati. \\
content & 1..* & contenuto il cui trasferimento e' governato. \\
\bottomrule
\end{tabularx}

\subsection{Invarianti e confini}
\begin{itemize}
  \item \texttt{transfer != produce}: produrre un risultato non prova che esso venga trasferito a un destinatario.
  \item \texttt{transfer != observe}: leggere qualcosa non prova una conveyance source-to-destination.
  \item \texttt{transfer != consumeService}: utilizzare un service non coincide con trasferire content.
  \item \texttt{transfer != initiate}: avviare/invocare qualcosa non implica conveyance.
  \item interaction mode, channel, retries, timeout, queueing, acknowledgement e failure semantics si aggiungono solo se governati e rappresentabili nel livello corretto.
\end{itemize}

\subsection{Forbidden inference}
Non inferire protocollo, sincronicita', cifratura, affidabilita', successo, persistenza, ownership della source/destination o topologia intermedia.

% -----------------------------------------------------------------------------
\section{Operator: produce}
\stabletag\quad\Rtwentyfivetag\quad\pendingexampletag

\subsection{Definizione}
\texttt{produce} rappresenta un actor/capability che rende disponibile un risultato governato. Il risultato puo' gia' avere identita' semantica senza essere ``creato'' nel senso di nuova occurrence/item project-semantic.

\subsection{Firma canonica}
\begin{lstlisting}
produce
    actor  -> BAReferent [1]
    result -> BAReferent [1..*]
    input  -> BAReferent [0..*]
\end{lstlisting}

\subsection{Campi}
\begin{tabularx}{\textwidth}{L{36mm}L{24mm}Y}
\toprule
actor & 1 & soggetto/capability che produce il risultato. \\
result & 1..* & output/result reso disponibile. \\
input & 0..* & input governati effettivamente usati, non semplicemente disponibili. \\
\bottomrule
\end{tabularx}

\subsection{Invarianti e confini}
\begin{itemize}
  \item \texttt{produce != create}: output availability non prova nuova identity/occurrence.
  \item \texttt{produce != transfer}: la disponibilita' di un result non identifica source/destination di conveyance.
  \item \texttt{produce != decisionRule}: un risultato prodotto non implica un mapping condizionale governato.
  \item \texttt{produce != observe}: un result di una query puo' essere prodotto, ma la lettura dell'esistente e il result prodotto sono semantiche separabili.
\end{itemize}

\subsection{Forbidden inference}
Non inferire algoritmo, score, ranking, persistenza, destinazione, nuova identity o condition logic non esplicitamente governati.

% -----------------------------------------------------------------------------
\section{Operator: create}
\stabletag\quad\Rtwentyfivetag\quad\pendingexampletag

\subsection{Definizione}
\texttt{create} rappresenta l'istituzione di una nuova identita'/occurrence/item project-semantic che prima dell'azione non esisteva nel significato governato.

\subsection{Firma canonica}
\begin{lstlisting}
create
    actor   -> BAReferent [1]
    created -> BAReferent [1..*]
\end{lstlisting}

\subsection{Campi}
\begin{tabularx}{\textwidth}{L{36mm}L{24mm}Y}
\toprule
actor & 1 & soggetto/capability che stabilisce la nuova identity/occurrence. \\
created & 1..* & item/event project-semantic creato. \\
\bottomrule
\end{tabularx}

\subsection{Invarianti e confini}
\begin{itemize}
  \item \texttt{create != produce}: generare un output non dimostra nuova identity.
  \item \texttt{create != transition}: cambiare stato dello stesso subject non crea un nuovo subject.
  \item \texttt{create != storedIn}: persistenza/at-rest association non prova creazione.
  \item nomi tecnici come factory, generator, constructor o insert non sono sufficienti.
\end{itemize}

\subsection{Forbidden inference}
Non inferire ID generation, storage technology, lifecycle successivo, ownership o correlation con altri elementi se non governati.

% -----------------------------------------------------------------------------
\section{Operator: observe}
\stabletag\quad\Rtwentyfivetag\quad\pendingexampletag

\subsection{Definizione}
\texttt{observe} rappresenta read/query/inspection esplicita di significato esistente senza asserire che l'osservazione crei o modifichi l'oggetto osservato.

\subsection{Firma canonica}
\begin{lstlisting}
observe
    actor    -> BAReferent [1]
    observed -> BAReferent [1..*]
    result   -> BAReferent [0..*]
\end{lstlisting}

\subsection{Campi}
\begin{tabularx}{\textwidth}{L{36mm}L{24mm}Y}
\toprule
actor & 1 & soggetto/capability che legge o ispeziona. \\
observed & 1..* & significato esistente oggetto dell'osservazione. \\
result & 0..* & eventuale result distinto prodotto dall'osservazione; CMD-OP04 ne mantiene aperta la revisione. \\
\bottomrule
\end{tabularx}

\subsection{Invarianti e confini}
\begin{itemize}
  \item \texttt{observe != input}: il fatto che un dato sia input non prova un'operazione di lettura semanticamente governata.
  \item \texttt{observe != transfer}: ricevere content non equivale automaticamente a osservarlo.
  \item \texttt{observe != selection}: query/retrieval puo' avvenire dopo una selection; non stabilisce la membership selezionata.
  \item \texttt{observe != transition}: inspection non implica state change.
\end{itemize}

\subsection{Open delta}
\opentag\quad \texttt{CMD-OP04}: la rimozione di \texttt{observe.result} resta deferita a integrated review; questa candidate preserva quindi la firma R3.

% -----------------------------------------------------------------------------
\section{Operator: transition}
\stabletag\quad\Rtwentyfivetag\quad\pendingexampletag

\subsection{Definizione}
\texttt{transition} rappresenta un cambiamento governato di stato/lifecycle dello \emph{stesso} subject project-semantic.

\subsection{Firma canonica}
\begin{lstlisting}
transition
    subject   -> BAReferent [1]
    toState   -> BAReferent [1]
    actor     -> BAReferent [0..1]
    fromState -> BAReferent [0..1]
\end{lstlisting}

\subsection{Campi}
\begin{tabularx}{\textwidth}{L{36mm}L{24mm}Y}
\toprule
subject & 1 & identita' che persiste attraverso il cambiamento. \\
toState & 1 & stato risultante governato. \\
actor & 0..1 & soggetto che causa/realizza il cambiamento, solo se governato. \\
fromState & 0..1 & stato precedente governato, se noto. \\
\bottomrule
\end{tabularx}

\subsection{Invarianti e confini}
\begin{itemize}
  \item stessa identity prima/dopo e' requisito centrale;
  \item selection di una versione non e' automaticamente transition;
  \item qualification result non e' automaticamente state transition;
  \item modifica di una property non prova lifecycle/state identity;
  \item supporto al rollback non prova che una transition specifica sia governata o automatica.
\end{itemize}

\subsection{Open delta}
\opentag\quad \texttt{CMD-OP05}: la distinzione state identity vs typed value resta deferita; nessuna estensione della firma R3 viene ammessa qui.

% -----------------------------------------------------------------------------
\section{Operator: correlate}
\stabletag\quad\Rtwentyfivetag\quad\pendingexampletag

\subsection{Definizione}
\texttt{correlate} preserva un binding di same-context tra uno o piu' elementi e una identita' di contesto condivisa.

\subsection{Firma canonica}
\begin{lstlisting}
correlate
    correlatedItem     -> BAReferent [1..*]
    correlationContext -> BAReferent [1]
\end{lstlisting}

\subsection{Campi}
\begin{tabularx}{\textwidth}{L{43mm}L{24mm}Y}
\toprule
correlatedItem & 1..* & elemento/i che devono restare nello stesso contesto. \\
correlationContext & 1 & identita' che definisce il contesto condiviso. \\
\bottomrule
\end{tabularx}

\subsection{Invarianti e confini}
\begin{itemize}
  \item \texttt{correlate != reference}: il reference puo' essere direzionale senza same-context invariant.
  \item \texttt{correlate != dependOn}: same-context binding non implica prerequisite.
  \item \texttt{correlate != transfer}: binding non implica conveyance.
  \item non usare correlate come generic association fallback.
\end{itemize}

% -----------------------------------------------------------------------------
\section{Operator: reference}
\stabletag\quad\Rtwentyfivetag\quad\pendingexampletag

\subsection{Definizione}
\texttt{reference} rappresenta un riferimento direzionale esplicito da un referencer a uno o piu' target identificati, quando non e' governata una semantica piu' forte.

\subsection{Firma canonica}
\begin{lstlisting}
reference
    referencer -> BAReferent [1]
    referenced -> BAReferent [1..*]
\end{lstlisting}

\subsection{Campi}
\begin{tabularx}{\textwidth}{L{36mm}L{24mm}Y}
\toprule
referencer & 1 & elemento che mantiene il riferimento. \\
referenced & 1..* & target del riferimento. \\
\bottomrule
\end{tabularx}

\subsection{Invarianti e confini}
\begin{itemize}
  \item \texttt{reference != correlate}: non garantisce same-context binding.
  \item \texttt{reference != dependOn}: non garantisce prerequisite.
  \item \texttt{reference != assignResponsibility}: non assegna authority/responsibility.
  \item non usare reference come fallback per relazioni semanticamente non comprese.
\end{itemize}

% -----------------------------------------------------------------------------
\section{Operator: dependOn}
\stabletag\quad\Rtwentyfivetag\quad\pendingexampletag

\subsection{Definizione}
\texttt{dependOn} rappresenta una relazione di prerequisite: il dependent richiede il prerequisite per il significato governato in scope.

\subsection{Firma canonica}
\begin{lstlisting}
dependOn
    dependent    -> BAReferent [1]
    prerequisite -> BAReferent [1..*]
\end{lstlisting}

\subsection{Campi}
\begin{tabularx}{\textwidth}{L{36mm}L{24mm}Y}
\toprule
dependent & 1 & significato che necessita del prerequisite. \\
prerequisite & 1..* & prerequisite governati richiesti. \\
\bottomrule
\end{tabularx}

\subsection{Invarianti e confini}
\begin{itemize}
  \item temporal order alone non e' dependency;
  \item \texttt{precedes != dependOn};
  \item pipeline membership non e' dependency;
  \item \texttt{dependOn != consumeService}: un prerequisite puo' non essere un service consumato;
  \item la transitive closure puo' essere projection-derived e non deve essere materializzata come truth senza necessita'.
\end{itemize}

% -----------------------------------------------------------------------------
\section{Operator: consumeService}
\stabletag\quad\Rtwentyfivetag\quad\pendingexampletag

\subsection{Definizione}
\texttt{consumeService} rappresenta l'uso effettivo di un service/capability da parte di uno o piu' consumer, senza trasferire automaticamente ownership o provider responsibility.

\subsection{Firma canonica}
\begin{lstlisting}
consumeService
    consumer -> BAReferent [1..*]
    service  -> BAReferent [1]
    provider -> BAReferent [0..1]
\end{lstlisting}

\subsection{Campi}
\begin{tabularx}{\textwidth}{L{36mm}L{24mm}Y}
\toprule
consumer & 1..* & utilizzatore/i effettivi del service. \\
service & 1 & service/capability consumato. \\
provider & 0..1 & provider esplicitamente governato; non implicito. \\
\bottomrule
\end{tabularx}

\subsection{Invarianti e confini}
\begin{itemize}
  \item endpoint, componente, tecnologia, dataset o information item non sono automaticamente service;
  \item \texttt{consumeService != transfer};
  \item \texttt{consumeService != provideService candidate};
  \item consumption non implica ownership, administration o infrastructure responsibility;
  \item un consumer requirement non prova che il service garantisca la property richiesta.
\end{itemize}

% -----------------------------------------------------------------------------
\section{Operator: realize}
\stabletag\quad\Rtwentyfivetag\quad\pendingexampletag

\subsection{Definizione}
\texttt{realize} rappresenta la materializzazione concreta di un significato astratto che esiste indipendentemente come commitment/capability/contract governato.

\subsection{Firma canonica}
\begin{lstlisting}
realize
    abstract    -> BAReferent [1]
    realization -> BAReferent [1..*]
\end{lstlisting}

\subsection{Campi}
\begin{tabularx}{\textwidth}{L{36mm}L{24mm}Y}
\toprule
abstract & 1 & significato astratto indipendentemente governato. \\
realization & 1..* & elemento/i concreti che lo materializzano. \\
\bottomrule
\end{tabularx}

\subsection{Invarianti e confini}
\begin{itemize}
  \item \texttt{realize != perform/execute}: realizzare una capability non equivale a rappresentare l'atto di eseguirla;
  \item \texttt{realize != classify}: non assegna semplicemente un kind;
  \item \texttt{realize != assignResponsibility}: materialization non equivale a responsibility placement;
  \item deve esistere un abstract meaning indipendentemente governato.
\end{itemize}

% -----------------------------------------------------------------------------
\section{Operator: assignResponsibility}
\stabletag\quad\Rtwentyfivetag\quad\pendingexampletag

\subsection{Definizione}
\texttt{assignResponsibility} colloca una responsibility o authority governata su una o piu' responsible party per uno scope e un responsibility kind espliciti.

\subsection{Firma canonica}
\begin{lstlisting}
assignResponsibility
    responsibleParty    -> BAReferent [1..*]
    responsibilityScope -> BAReferent [1]
    responsibilityKind  -> BAReferent [1]
\end{lstlisting}

\subsection{Campi}
\begin{tabularx}{\textwidth}{L{45mm}L{24mm}Y}
\toprule
responsibleParty & 1..* & parte/i cui e' attribuita la responsibility. \\
responsibilityScope & 1 & ambito governato della responsibility. \\
responsibilityKind & 1 & tipo di responsibility/authority governato. \\
\bottomrule
\end{tabularx}

\subsection{Invarianti e confini}
\begin{itemize}
  \item provider != responsible party per default;
  \item performer/actor != responsible party per default;
  \item producer/consumer != responsible party per default;
  \item realization != responsibility;
  \item mera partecipazione non assegna authority.
\end{itemize}

% -----------------------------------------------------------------------------
\section{Operator: constrain}
\stabletag\quad\Rtwentyfivetag\quad\pendingexampletag

\subsection{Definizione}
\texttt{constrain} rappresenta una restriction governata, riusabile e interrogabile applicata a un target identificato.

\subsection{Firma canonica}
\begin{lstlisting}
constrain
    constraintTarget -> BAReferent [1]
    constraintValue  -> controlled value/structure [1..*]
\end{lstlisting}

\subsection{Campi}
\begin{tabularx}{\textwidth}{L{42mm}L{24mm}Y}
\toprule
constraintTarget & 1 & significato esatto su cui insiste la restriction. \\
constraintValue & 1..* & valore, dominio o struttura di constraint governata. \\
\bottomrule
\end{tabularx}

\subsection{Invarianti e confini}
\begin{itemize}
  \item \texttt{constrain != decisionCondition}: una condition locale non e' automaticamente reusable constraint;
  \item \texttt{constrain != decisionRule}: un dominio ammesso non definisce necessariamente come scegliere un risultato;
  \item selection bound puo' coesistere con constrain ma non assorbe selected membership;
  \item \texttt{NOT\_GOVERNED} non viene trasformato in constraint negativo.
\end{itemize}

% -----------------------------------------------------------------------------
\section{Operator: classify}
\stabletag\quad\Rtwentyfivetag\quad\pendingexampletag

\subsection{Definizione}
\texttt{classify} assegna a un referent uno o piu' semantic kind/category methodology-neutral governati.

\subsection{Firma canonica}
\begin{lstlisting}
classify
    classifiedReferent -> BAReferent [1]
    semanticKind       -> BAReferent [1..*]
\end{lstlisting}

\subsection{Campi}
\begin{tabularx}{\textwidth}{L{42mm}L{24mm}Y}
\toprule
classifiedReferent & 1 & elemento cui viene assegnato il kind. \\
semanticKind & 1..* & categoria/kind governati. \\
\bottomrule
\end{tabularx}

\subsection{Invarianti e confini}
\begin{itemize}
  \item \texttt{classify != selection}: assegnare un kind non sceglie un membro tra candidati;
  \item \texttt{classify != transition}: cambiare category non prova lifecycle state change;
  \item \texttt{classify != generic property assignment};
  \item il kind deve essere methodology-neutral e materialmente utile, non introdotto solo per il tool.
\end{itemize}

% -----------------------------------------------------------------------------
\section{Operator: decisionRule}
\stabletag\quad\Rtwentyfivetag\quad\pendingexampletag

\subsection{Definizione}
\texttt{decisionRule} rappresenta un mapping governato da input/conditions a un risultato governato. E' appropriato quando il progetto governa non soltanto il dominio dei risultati, ma la regola che determina il risultato.

\subsection{Firma canonica}
\begin{lstlisting}
decisionRule
    actor  -> BAReferent [1]
    input  -> BAReferent [1..*]
    result -> BAReferent [1]

    rule
        IF    decisionCondition
        THEN  resultAssignment [1..*]
        ELSE  resultAssignment [0..*]
\end{lstlisting}

\subsection{Campi top-level}
\begin{tabularx}{\textwidth}{L{36mm}L{24mm}Y}
\toprule
actor & 1 & soggetto/capability che applica la regola. \\
input & 1..* & input governati effettivamente usati dal mapping. \\
result & 1 & risultato governato dal rule mapping. \\
rule & 1 local structure & struttura operator-local di condition e assignment. \\
\bottomrule
\end{tabularx}

\subsection{Invarianti e confini}
\begin{itemize}
  \item conditional action non e' automaticamente decisionRule;
  \item non-sufficiency statement non e' automaticamente decisionRule;
  \item ranking/selection non e' automaticamente decisionRule;
  \item capability/support statement non e' automaticamente decisionRule;
  \item un allowed result domain da solo non definisce il mapping;
  \item property/value interni non vengono inventati per soddisfare la sintassi.
\end{itemize}

% -----------------------------------------------------------------------------
\section{Condition language di decisionRule}
\stabletag\quad\Rtwentyfivetag

Le condition form sono strutture locali. Non aumentano il conteggio dei top-level operator.

\subsection{comparison}
\begin{lstlisting}
comparison
    referent      -> BAReferent [1]
    property      -> controlled semantic key [1]
    comparisonKey -> equals | notEquals
    value         -> typed local value | BAReferent [1]
\end{lstlisting}

\begin{tabularx}{\textwidth}{L{38mm}Y}
\toprule
referent & elemento la cui property viene confrontata. \\
property & semantic key controllata; resta obbligatoria nel contratto corrente. \\
comparisonKey & operazione di confronto ammessa. Current: \texttt{equals}, \texttt{notEquals}. \\
value & valore tipizzato o BAReferent contro cui avviene il confronto. \\
\bottomrule
\end{tabularx}

\begin{ddtaopen}[title={CL-01 - ordered/scalar comparison vocabulary}]
Il ciclo R25 ha prodotto forte evidence per confronti ordinati/scalari che superano \texttt{equals/notEquals}. Le relazioni concettuali \texttt{<}, \texttt{>}, \texttt{<=}, \texttt{>=} sono pressure reale, ma il controlled vocabulary normativo e la relativa typing discipline non sono ancora chiusi. Non usarli come vocabulary BA2 current senza formal delta.
\end{ddtaopen}

\subsection{satisfies}
\stabletag
\begin{lstlisting}
satisfies
    subject   -> BAReferent [1]
    condition -> BAReferent [1]
\end{lstlisting}
Serve quando la fonte governa che un subject soddisfi una condition identificata, ma non governa un modello property/value interno che autorizzi una \texttt{comparison}.

\subsection{allOf}
\stabletag\quad\Rtwentyfivetag
\begin{lstlisting}
allOf
    childCondition -> decisionCondition [2..*]
\end{lstlisting}
Tutte le condizioni figlie devono valere. R25 fornisce forte evidence di conjunction semantics, senza trasformare \texttt{allOf} in top-level operator.

\subsection{anyOf}
\stabletag
\begin{lstlisting}
anyOf
    childCondition -> decisionCondition [2..*]
\end{lstlisting}
Almeno una condizione figlia deve valere. Il costrutto e' gia' corrente, ma R25 lo mantiene in stato \texttt{PENDING}: il positive control e' ancora insufficiente. Questo non equivale a rimozione dal contratto corrente.

\subsection{not}
\stabletag
\begin{lstlisting}
not
    childCondition -> decisionCondition [1]
\end{lstlisting}
Negazione locale di una condition. La frase ``A non e' sufficiente per B'' non equivale a \texttt{not(B)} e non deve essere compressa in una negazione logica falsa.

\subsection{resultAssignment}
\stabletag
\begin{lstlisting}
resultAssignment
    target -> BAReferent [1]
    value  -> typed local value | BAReferent [1]
\end{lstlisting}
Rappresenta l'assegnazione del risultato governato nel ramo THEN/ELSE. Non autorizza property assignment generico fuori da \texttt{decisionRule}.

% -----------------------------------------------------------------------------
\section{Struttura locale candidate: selection}
\candidatetag\quad\Rtwentyfivetag\quad\pendingexampletag

R25 ha dimostrato che bounded/ranked/recency selection contiene semantica non eliminabile in diversi controlli, senza giustificare un nuovo top-level operator.

\subsection{Working structure}
\begin{lstlisting}
selection
    candidatePopulation -> OPEN TYPE [required]
    selectionBasis      -> OPEN TYPE [required]
    selectedResult      -> OPEN TYPE [required]
    bound               -> OPEN TYPE [optional, only if governed]
    orderingBasis       -> OPEN TYPE [optional, only if governed]
\end{lstlisting}

\begin{ddtaopen}[title={Typing non ancora congelato}]
Questa struttura e' \texttt{REUSABLE\_SELECTION\_STRUCTURE / NON-NORMATIVE}. R25 non decide ancora se i campi siano BAReferent, local structured values, condition-language structures o altro tipo gia' governato. \texttt{criterion} e \texttt{selectionBasis} restano una naming decision da chiudere nell'integrazione metodologica.
\end{ddtaopen}

\subsection{Distinzioni obbligatorie}
\begin{lstlisting}
ranking        != selection
bound          != selection
retrieval      != selection
filtering      != selection
best selection != full ranking
\end{lstlisting}

La membership selezionata deve restare distinguibile dal criterio, dall'eventuale ordinamento e dall'eventuale bound.

% -----------------------------------------------------------------------------
\section{Pressure aperta: pipeline composition e order}
\opentag\quad\Rtwentyfivetag\quad\pendingexampletag

\texttt{PR-02 PIPELINE\_COMPOSITION\_ORDER} resta aperta. La guida deve quindi preservare le distinzioni senza introdurre prematuramente \texttt{pipeline} o \texttt{precedes} come top-level operator.

\begin{lstlisting}
pipeline membership
    != pipeline ordering
    != precedes
    != dependOn
    != transfer
    != initiate
\end{lstlisting}

\subsection{Domande da porre}
\begin{itemize}
  \item La fonte governa soltanto appartenenza a una stessa composizione?
  \item Governa un ordine totale o parziale?
  \item L'ordine e' semantico o soltanto descrittivo/implementativo?
  \item B richiede A come prerequisite, oppure A e' soltanto precedente?
  \item Esiste conveyance tra stadi oppure solo successione logica?
  \item La composizione richiede una reusable local structure oppure e' proiettabile da proposition gia' accettate?
\end{itemize}

\begin{ddtaopen}[title={Current disposition}]
Nuovo top-level operator: \texttt{NOT JUSTIFIED}. Reusable local composition structure: possibile, ma non ancora chiusa. Test successivo: membership vs ordering; \texttt{precedes} vs \texttt{dependOn}.
\end{ddtaopen}

% -----------------------------------------------------------------------------
\section{Pressure e strutture trasversali ancora aperte}

\begin{longtable}{L{25mm}L{52mm}L{62mm}}
\toprule
\textbf{ID} & \textbf{Topic} & \textbf{Stato / disciplina} \\
\midrule\endhead
PR-01 & function/process/behavior identity / \texttt{perform?} & PRESSURE RETAINED; execution evidence esiste ma irreducibility non dimostrata. \\
PR-04 & boundary/interaction association & OPEN; distinguere crossing, transfer, interface e initiation. \\
PR-07 & structured information contract & OPEN; puo' appartenere a structure/schema/composition anziche' dynamic relation. \\
PR-09 & acquisition/refresh semantics & OPEN; esaurire operatori correnti prima di introdurre primitive. \\
PR-12 & negative implication/non-sufficiency & OPEN; non comprimere in \texttt{not}. \\
PR-14 & applicability/configuration binding & OPEN; provare qualifier/reusable structure prima di nuovo operatore. \\
OBS-OT-01 & operation target / effect scope & LEVEL UNRESOLVED; verificare local structure/composition. \\
\bottomrule
\end{longtable}

% -----------------------------------------------------------------------------
\section{Candidate construct R25 non ammessi}

\subsection{CC-01 consumeData}
\rejectedtag

Disposition: \texttt{REJECT\_REDUNDANT}. Nei casi correnti il non-destructive data/evidence use e' rappresentabile senza una nuova primitive top-level; destructive-consume semantics non risultano governate.

\subsection{CC-02 provideService}
\candidatetag

Positive provider/service evidence sopravvive l'exhaustion corrente, ma formal gate G1--G8, compatibility e final signature sono pendenti.

\begin{ddtaopen}[title={Signature non congelata}]
Questa guida non inventa roleKey o cardinalita' di \texttt{provideService}. Fino alla formal admission, il candidate resta evidence e non vocabulary BA2 corrente.
\end{ddtaopen}

\subsection{CC-03 storedIn}
\candidatetag

L'associazione at-rest sopravvive l'exhaustion corrente, ma formal admission, compatibility e final signature restano pendenti.

\begin{ddtaopen}[title={Signature non congelata}]
Questa guida non inventa un modello container/item, storage owner, persistence semantics o lifecycle. \texttt{storedIn} resta \texttt{TESTED\_POSITIVE / NOT\_ADMITTED}.
\end{ddtaopen}

\subsection{CC-04 initiate}
\candidatetag

\subsubsection{Working minimum signature}
\begin{lstlisting}
initiate
    initiator -> BAReferent [1]
    initiated -> BAReferent [1]
\end{lstlisting}

\subsubsection{Semantica minima}
Rappresenta una relazione diretta in cui un referent avvia/richiede l'avvio di un altro significato governato. La signature non distingue ancora request vs start come kind separati.

\subsubsection{Non implica}
\begin{lstlisting}
success
completion
transfer
service consumption
create
transition
automatic trigger
synchronous execution
\end{lstlisting}

Disposition: \texttt{TESTED\_POSITIVE / NOT\_ADMITTED}; formal G1--G8 e cross-corpus compatibility sono ancora richiesti.

% -----------------------------------------------------------------------------
\section{Working hypothesis respinte}
\rejectedtag

\begin{longtable}{L{45mm}L{95mm}}
\toprule
\textbf{Hypothesis} & \textbf{Disposition corrente} \\
\midrule\endhead
generic \texttt{consume} & NOT JUSTIFIED. \\
generic \texttt{modify} & REJECTED - false semantics come generic operator. \\
generic \texttt{persist} & NOT JUSTIFIED AS GENERIC OPERATOR. \\
generic \texttt{expose} & NOT JUSTIFIED - current corpus decomposable. \\
generic \texttt{trigger} & NOT JUSTIFIED AS SEPARATE PRIMITIVE; directed activation pressure routed to candidate \texttt{initiate} + condition. \\
verb-specific \texttt{invoke} & ABSORBED BY CC-04 \texttt{initiate}. \\
\texttt{initiationKind REQUEST|START} & REJECT REDUNDANT / OVER-SPECIFIC nel current evidence set. \\
\bottomrule
\end{longtable}

% -----------------------------------------------------------------------------
\section{BA3 - provenance, origin state e revalidation}
\stabletag

Ogni BA element meaning-bearing deve risolvere a una baseline governata e alla provenance richiesta dal contratto BA3.

\subsection{Origin states}
\begin{lstlisting}
GROUNDED
DERIVED
DIAGNOSTIC_UNRESOLVED
\end{lstlisting}

\subsection{Tre responsabilita' distinte}
\begin{lstlisting}
sourceLink
    dove e' supportato il significato

derivationBasisBinding
    quali elementi governati e quali ruoli fondano una derivazione

revalidationContext
    quale change/review context determina freshness e revalidation
\end{lstlisting}

Provenance non deve simulare una relazione BA2 mancante.

\subsection{Derived semantics}
Un elemento \texttt{DERIVED} richiede almeno:
\begin{itemize}
  \item derivation rule identificabile e versionata;
  \item basis role-bound;
  \item applicability contract;
  \item output semantic contract;
  \item lineage fino a governed source.
\end{itemize}

\subsection{Change e revalidation}
\begin{lstlisting}
source change
    -> explicit impact / revalidation context
    -> impacted BA elements
    -> PENDING_REVIEW / STALE where applicable
    -> RETAIN / REPLACE / RETIRE
    -> accepted BA for new baseline
    -> rebuild projections
\end{lstlisting}

Una accepted BA storica non viene mutata silenziosamente come se le analisi precedenti avessero usato la nuova semantica.

% -----------------------------------------------------------------------------
\section{Diagnostics e unresolved semantics}
\stabletag

Vocabolario operativo:
\begin{lstlisting}
MULTIPLE MATERIAL BA CANDIDATES
BA REQUIRES UNSUPPORTED FACT
BA LOSES GOVERNED MATERIAL DISTINCTION
BA EXPOSES SOURCE CONFLICT
NOT SPECIFIED
\end{lstlisting}

\begin{ddtacore}[title={NOT SPECIFIED}]
\texttt{NOT SPECIFIED} non significa false, absent o denied. Significa che il significato necessario non e' determinato dalla baseline governata corrente.
\end{ddtacore}

\subsection{Quando preferire una diagnosis a una proposition}
Usare una diagnosis quando la source non supporta una scelta univoca tra interpretazioni materialmente differenti, quando la firma corrente richiederebbe un fatto non governato o quando il metodo perderebbe una distinzione reale.

% -----------------------------------------------------------------------------
\section{BA5 - authoring canonico}
\stabletag

I semantic key devono essere canonici nel proprio semantic domain.

\begin{itemize}
  \item project referent names;
  \item BA2 operator keys;
  \item operator-scoped role keys;
  \item semantic kind keys;
  \item controlled property keys;
  \item controlled value domains;
  \item local-structure keys quando formalmente ammessi.
\end{itemize}

Sinonimia, traduzione, naming similarity o LLM similarity non stabiliscono equivalenza normativa automaticamente.

\begin{ddtacore}[title={Registry boundary}]
BA5 governa canonicalita' e resolution; non puo' ridefinire semantica, roleKey o cardinalita' di un operatore BA2. Per questi aspetti l'authority corrente resta BA2 R3.
\end{ddtacore}

% -----------------------------------------------------------------------------
\section{BA4 - projection boundary}
\stabletag

\begin{lstlisting}
governed documentation
    -> accepted BA
        -> projection
            -> method-owned interpretation
\end{lstlisting}

Ogni meaning-bearing projection item deve tracciare alla BA.

\subsection{Coverage modes}
\begin{lstlisting}
EXHAUSTIVE_FOR_DECLARED_SCOPE
SELECTIVE
\end{lstlisting}

Una selective projection puo' omettere elementi senza dichiararli falsi o assenti. Una method-owned interpretation puo' aggiungere semantica del metodo downstream, ma deve restare trace-bound e non diventare BA/project authority.

\subsection{Projection non autorizza inference upstream}
Un DFD, un graph, una threat taxonomy o una vista architetturale non possono costringere la BA ad aggiungere endpoint, trust boundary, actor, data flow o service relationship non governati.

% -----------------------------------------------------------------------------
\section{Regression e acceptance della Base Analysis}
\stabletag

Dopo una material BA change o contract refinement:
\begin{enumerate}
  \item ricostruire o aggiornare la BA contro la source pin;
  \item verificare ogni proposition materiale;
  \item verificare diagnostics e forbidden inference;
  \item controllare property/constraint scope;
  \item controllare request/case/identity/evaluation bindings;
  \item verificare signatures e cardinalita';
  \item rieseguire usefulness/projection checks appropriati;
  \item cercare counterexample e semantic loss;
  \item riaprire soltanto il minimo owning layer se necessario.
\end{enumerate}

\subsection{Checklist pre-acceptance}
\begin{enumerate}
  \item authority e source revision sono pinned?
  \item ogni referent ha identity solo quando necessaria?
  \item ogni proposition usa operator/role canonici?
  \item tutte le cardinalita' required sono source-grounded?
  \item ruoli optional sono presenti solo se governati?
  \item e' stato introdotto un fatto non governato?
  \item un \texttt{NOT SPECIFIED} e' diventato default?
  \item provenance e revalidation sono sufficienti?
  \item local structure candidate e' stata scambiata per normative vocabulary?
  \item una candidate R25 e' stata presentata come admitted?
  \item una projection puo' consumare la BA senza reinterpretare prose per shared meaning gia' rappresentato?
  \item esiste un counterexample che richiede targeted reopen?
\end{enumerate}

% -----------------------------------------------------------------------------
\section{Open register R25 da preservare nella prossima revisione}
\opentag

\begin{longtable}{L{25mm}L{55mm}L{58mm}}
\toprule
\textbf{ID} & \textbf{Tema} & \textbf{Stato corrente} \\
\midrule\endhead
PR-01 & perform / execution identity & PRESSURE RETAINED; irreducibility not demonstrated. \\
PR-02 & pipeline composition/order & OPEN. \\
PR-03 & interface/path/invocation & RECONCILED / CC-04. \\
PR-04 & boundary/interaction & OPEN. \\
PR-05 & ordered comparison & ROUTE CL-01. \\
PR-06 & scalar/property comparison & ROUTE CL-01. \\
PR-07 & structured information contract & OPEN. \\
PR-08 & at-rest storage & RECONCILED / CC-03. \\
PR-09 & acquisition/refresh & OPEN. \\
PR-10 & conditional action trigger & RECONCILED / CC-04 + CONDITION. \\
PR-11 & data/evidence consumption & RECONCILED. \\
PR-12 & negative implication/non-sufficiency & OPEN. \\
PR-13 & bounded/ranked/recency selection & REUSABLE SELECTION STRUCTURE / NON-NORMATIVE. \\
PR-14 & applicability/configuration binding & OPEN. \\
CMD-OP04 & remove observe.result & DEFERRED. \\
CMD-OP05 & transition state/value refinement & DEFERRED. \\
OBS-OT-01 & operation-target/effect-scope & LEVEL UNRESOLVED. \\
\bottomrule
\end{longtable}

% -----------------------------------------------------------------------------
\section{Cosa questa revisione non chiude}

Questa candidate non chiude:
\begin{itemize}
  \item BA6 integrated acceptance;
  \item formal admission di \texttt{provideService}, \texttt{storedIn}, \texttt{initiate};
  \item ordered/scalar comparator vocabulary oltre \texttt{equals/notEquals};
  \item typing finale della selection structure;
  \item pipeline membership/order representation;
  \item function/process/behavior identity pressure;
  \item boundary/interactions, structured information, acquisition/refresh, negative non-sufficiency, applicability binding;
  \item CMD-OP04, CMD-OP05 e OBS-OT-01.
\end{itemize}

\section{Gate per la prossima revisione con esempi}
\pendingexampletag

La prossima revisione deve mantenere integralmente questa struttura e aggiungere, senza rimuovere semantica:
\begin{enumerate}
  \item almeno un positive o boundary control per ciascun top-level operator;
  \item Facial Access come regression case storico dove disponibile;
  \item DermaTriage come secondo/control case dove disponibile;
  \item source clause $\rightarrow$ semantic facts $\rightarrow$ BAReferent $\rightarrow$ BAProposition;
  \item rejected alternative routing;
  \item forbidden inference;
  \item errori frequenti;
  \item esempi specifici per condition language, selection e pipeline/order;
  \item dichiarazione esplicita ``no clean positive control'' quando il corpus non ne offre uno.
\end{enumerate}

\begin{ddtacore}[title={Acceptance criterion della Complete Guide companion}]
La guida sara' candidata alla promotion soltanto quando un lettore potra' sia ricostruire il contratto BA completo dalla guida, sia imparare ad applicarlo correttamente senza dover consultare conversazioni o working findings non integrati.
\end{ddtacore}

\appendix
\section{Quick syntax sheet - top-level operators}

\begin{lstlisting}
transfer
    behavior    -> BAReferent [0..1]
    source      -> BAReferent [1]
    destination -> BAReferent [1..*]
    content     -> BAReferent [1..*]

produce
    actor  -> BAReferent [1]
    result -> BAReferent [1..*]
    input  -> BAReferent [0..*]

create
    actor   -> BAReferent [1]
    created -> BAReferent [1..*]

observe
    actor    -> BAReferent [1]
    observed -> BAReferent [1..*]
    result   -> BAReferent [0..*]

transition
    subject   -> BAReferent [1]
    toState   -> BAReferent [1]
    actor     -> BAReferent [0..1]
    fromState -> BAReferent [0..1]

correlate
    correlatedItem     -> BAReferent [1..*]
    correlationContext -> BAReferent [1]

reference
    referencer -> BAReferent [1]
    referenced -> BAReferent [1..*]

dependOn
    dependent    -> BAReferent [1]
    prerequisite -> BAReferent [1..*]

consumeService
    consumer -> BAReferent [1..*]
    service  -> BAReferent [1]
    provider -> BAReferent [0..1]

realize
    abstract    -> BAReferent [1]
    realization -> BAReferent [1..*]

assignResponsibility
    responsibleParty    -> BAReferent [1..*]
    responsibilityScope -> BAReferent [1]
    responsibilityKind  -> BAReferent [1]

constrain
    constraintTarget -> BAReferent [1]
    constraintValue  -> controlled value/structure [1..*]

classify
    classifiedReferent -> BAReferent [1]
    semanticKind       -> BAReferent [1..*]

decisionRule
    actor  -> BAReferent [1]
    input  -> BAReferent [1..*]
    result -> BAReferent [1]
    rule   -> operator-local structure
\end{lstlisting}

\section{Quick syntax sheet - condition language}
\begin{lstlisting}
comparison
    referent      -> BAReferent [1]
    property      -> controlled semantic key [1]
    comparisonKey -> equals | notEquals
    value         -> typed local value | BAReferent [1]

satisfies
    subject   -> BAReferent [1]
    condition -> BAReferent [1]

allOf
    childCondition -> decisionCondition [2..*]

anyOf
    childCondition -> decisionCondition [2..*]

not
    childCondition -> decisionCondition [1]

resultAssignment
    target -> BAReferent [1]
    value  -> typed local value | BAReferent [1]
\end{lstlisting}

\section{Quick status sheet - non-current constructs}
\begin{lstlisting}
selection
    REUSABLE_SELECTION_STRUCTURE / NON-NORMATIVE

provideService
    TESTED_POSITIVE / NOT_ADMITTED
    signature = NOT FROZEN

storedIn
    TESTED_POSITIVE / NOT_ADMITTED
    signature = NOT FROZEN

initiate
    TESTED_POSITIVE / NOT_ADMITTED
    working minimum:
        initiator -> BAReferent [1]
        initiated -> BAReferent [1]

pipeline/order
    OPEN
    no new top-level operator admitted
\end{lstlisting}

\end{document}
